% =====================================================================
%  Section 8 - Discussion and threats to validity (author's version,
%  facts checked against Sects. 4-7 and NoiseStudy/report.txt)
% =====================================================================
\section{Discussion and Threats to Validity}
\label{sec:discussion}

\paragraph{Why generational replacement fails.} Generational replacement replaces all the parents, meaning no reel set survives unchanged. Near the end of a run, when the reel sets are close to the targets, the children are mostly worse than their parents, and a higher mutation strength has a more severe impact. Unlike generational replacement, steady-state and elitist replacement never lose the best reel set, so a bad child costs just one evaluation, nothing more.

\paragraph{Selection only matters when replacement does not keep the best.} With steady-state and elitist replacement the best reel sets are always kept, but with generational replacement, selection is the only way good reel strips reach the next generation. When 10 parents are drawn from 10 members, Tournament ($k = 3$) picks the best member 3 times on average and misses it in $0.7^{10} \approx 3\%$ of generations. SUS picks it exactly twice. Roulette also picks it twice on average, but misses it in $0.8^{10} \approx 11\%$ of generations.
\paragraph{In the base stage, local search is enough.} Hill climbing is a steady-state GA with one member and no crossover (Sect.~\ref{sec:baselines}), so the only way the GA can beat it is if crossover or the population helps, but in the base stage neither did. Line-win RTP depends only on how many of each symbol a reel strip holds (Sect.~\ref{sec:noise}), and one mutation changes two of these counts by 1 in 50. So even without crossover, a chain of small mutation steps lowers the error. We think crossover adds very little near the end, but we never measured diversity or ran a GA without crossover.

\paragraph{In the free stage, a population and annealing help.} The free-game target stats conflict with each other (Sect.~\ref{sec:res-free}), so near the end every small change likely makes at least one of the stats worse. Since hill climbing only accepts a child that is not worse, it stays there. At $m = 3$ all six steady-state and elitist configurations beat it. Simulated annealing also works on a single reel set, but it accepts some worse children, and no GA configuration beats it at any $m$. So the gain seems to come from keeping or accepting reel sets other than the current best, which a population and annealing both do, and not from crossover.

\paragraph{Advice for slot mathematicians.} Start with a steady-state GA, an elitist GA or simulated annealing; all three reach a base error of about 0.1, e.g.\ a base RTP 1\% off its target. Do not use generational replacement without elitism. Changing 3 to 5 of the 250 symbols is a safe mutation strength for both stages. Also, re-simulate the final reel set with more spins on a new seed. Finally, if every good method stops at about the same error, as in our free stage (best medians 1.17 to 1.30), check the targets before changing the optimiser.

\paragraph{Threats to validity.} \emph{One game, one budget, one population size.} We tuned a simple line-win game against three stats per stage. More complex games may favour other methods. Every run had 1000 evaluations and $\mu = 10$. A larger budget might let crossover pay off, and a larger population gives more diversity but fewer generations; we tested neither. \emph{Untuned methods.} The GA used textbook defaults, and the annealing temperatures (0.14 to 0.002) were set once and never changed. Tuning could change the order among the tied methods, so our claim is about sensible defaults. \emph{Best $m$ chosen on the same seeds.} When each method is compared at its own best $m$, the same 30 seeds choose $m$ and measure the result. This favours the GA, whose best value is the lowest of 36 configurations (9 $\times$ 4 values of $m$) against 4 for simulated annealing. The base-stage draw holds despite this, and at a fixed $m = 5$ it is also a draw. \emph{Weighted-sum error and conflicting free targets.} The weights 10, 6 and 3 are our choice, and other weights could rank reel sets and methods differently. Only 16 of 2310 final free reel sets reach an error of 0.25 or less, so the free stage compares how methods trade one target against another. We did not prove that the free targets cannot all be met; the exact calculation of Sect.~\ref{sec:noise} could settle this. \emph{Noise floor.} Its standard deviations (0.034 and 0.10) come from a measurement made before the main experiment, and are above every value in the final noise study.